\section{The fixed map and its complete fibers}\label{sec:fixed-map}

The construction is easiest to understand backwards.  We want the inverse
problem over a target to be a single quintic equation in a marked root
\(S\), and we want the parameter \(a\) to enter only through the
coefficients of \(S^2\) and \(1\).  The quadratic gauge was designed
precisely so that these two coefficients are target coordinates, while the
derivative of the inverse equation cancels out of the Jacobian.

The centered identity from the introduction determines the seed.  Put
\[
 G(S)=S^5-\frac32S^4+\frac32S^3-\frac54S^2+\frac9{16}S.
 \tag{2.1}\label{eq:seed}
\]
The coefficients may look tailored, because they are: \(G\) is the fixed
part left after translating \(X^3(X^2+X+1)\) by \(X=S-\tfrac12\).

We first write the natural determinant-minus-two gauge map, and then make a
small linear normalization to obtain Jacobian one.  For source coordinates
\((x,y,z)\), set
\[
 t=1+xy,\qquad
 q=t^2z+\frac38y^2(1+3t),
 \tag{2.2}\label{eq:tq}
\]
and define \(F_0=(\Pi,B,C)\) by
\begin{align}
 \Pi={}&tq,\notag\\
 B={}&y+8xq-\frac{40}{9}tq
       -\frac{32}{3}t^2x^2q^4
       +\frac{80}{9}t^2x^3q^5,\notag\\
 C={}&x(5-3t)-\frac83x^3z
       +\frac{16}{3}(xq)^4-\frac{16}{3}(xq)^5.
 \tag{2.3}\label{eq:base-map}
\end{align}
These are polynomial coordinates over \(\Q\).
The formula is long, but there is no search hidden inside it: it is the
coefficientwise polynomial pullback of the one-variable inverse equation
below.

The determinant \(-2\) is the natural normalization of the inverse
equation.  It is not an arithmetic restriction.  Let
\[
 A(x,y,z)=\left(\frac{x}{2},y,z\right),
 \qquad
 L(\Pi,B,C)=(-\Pi,B,C),
 \tag{2.4}\label{eq:linear-normalization}
\]
and take as our fixed map
\[
 \boxed{\Phi=L\circ F_0\circ A.}
 \tag{2.5}\label{eq:normalized-map}
\]
Equations \eqref{eq:tq}--\eqref{eq:normalized-map} are an explicit polynomial formula for
\(\Phi\).  The normalization has been chosen so that it does not increase
the height of the target family.

\subsection{Inverse equation}

The inverse problem for \eqref{eq:base-map} is one-dimensional.  For an independent
target coordinate \(\Pi\), define
\[
 G_\Pi(S)=\frac9{16}S
 +\Pi\left(-\frac54S^2+\frac32S^3\right)
 -\frac32\Pi^4S^4+\Pi^5S^5
 \tag{2.6}\label{eq:g-pi}
\]
and
\[
 E_{\Pi,B,C}(S)
 =G_\Pi(S)-\frac9{32}(BS^2+C).
 \tag{2.7}\label{eq:inverse-equation}
\]

\begin{proposition}[Quadratic-gauge specialization]\label{prop:gauge}
The map \(F_0\) satisfies
\[
 \det DF_0=-2,\qquad \gdeg(F_0)=5.
 \tag{2.8}\label{eq:base-invariants}
\]
If \(\pi\ne0\) and \(E_{\pi,b,c}\) is squarefree, then there is a natural
fiber isomorphism
\[
 F_0^{-1}(\pi,b,c)
 \simeq\Spec\Q[S]/(E_{\pi,b,c}).
 \tag{2.9}\label{eq:base-fiber}
\]
\end{proposition}

\begin{proof}
We recall the short inverse-coordinate calculation, both to make the fiber
claim self-contained and to account for every source point.  On \(t\ne0\)
put
\[
 S=\frac{x}{t},\qquad Q=y+xq,\qquad
 D=1-SQ+\Pi S^2.
 \tag{2.10}\label{eq:chart}
\]
Since \(\Pi=tq\), one has \(D=t^{-1}\).  With
\[
 \beta(\Pi,S)
 =\frac{G_\Pi'(S)/(9/16)-1-\Pi S^2}{S},
 \tag{2.11}\label{eq:beta}
\]
direct expansion of \eqref{eq:base-map} gives
\[
 B=Q+\beta(\Pi,S),\qquad
 C=\frac{32}{9}G_\Pi(S)-BS^2.
 \tag{2.12}\label{eq:compact-target}
\]
The second identity is precisely \(E_{\Pi,B,C}(S)=0\), and differentiation
gives
\[
 D=\frac{16}{9}E_{\Pi,B,C}'(S).
 \tag{2.13}\label{eq:derivative}
\]
At fixed \(\Pi\), the Jacobian of \((S,Q)\mapsto(B,C)\) is \(-2D\);
the Jacobian of the source chart
\((x,y,z)\mapsto(\Pi,S,Q)\) is \(D^{-1}\).  Their product is \(-2\), proving
the first identity in \eqref{eq:base-invariants}.

If a source point lies over \((\pi,b,c)\) with \(\pi\ne0\), then
\(tq=\pi\), so \(t\) and \(q\) are units and the chart above contains the
entire fiber.  Conversely, let \(s\) be a root of \(E_{\pi,b,c}\) over a
commutative \(\Q\)-algebra and suppose \(E_{\pi,b,c}'(s)\) is a unit.  Put
\[
 \begin{gathered}
 d=\frac{16}{9}E_{\pi,b,c}'(s),\qquad
 Q=b-\beta(\pi,s),\qquad t=d^{-1},\\
 x=sd^{-1},\qquad y=Q-\pi s,\qquad q=\pi d,\\
 z=d^2\left(q-\frac38y^2(1+3t)\right).
 \end{gathered}
 \tag{2.14}\label{eq:reconstruction}
\]
Equations \eqref{eq:chart}--\eqref{eq:derivative} show directly that this reconstructs a
unique source point.  If \(E_{\pi,b,c}\) is squarefree, B\'ezout's identity
for \(E\) and \(E'\) makes \(E'(s)\) a unit over every test algebra.  Hence
\eqref{eq:reconstruction} is a natural inverse and proves \eqref{eq:base-fiber}.

Finally, the generic polynomial \eqref{eq:inverse-equation} has degree five in \(S\) and is
irreducible over \(\Q(\Pi,B,C)\).  Indeed, over
\(\Q(\Pi,B)[C]\) it is primitive and linear in \(C\); in a nontrivial
factorization one factor would be independent of \(C\), and comparison of
the coefficient of \(C\) would make that factor divide a nonzero scalar.
The reconstruction identifies the actual source field with the field
obtained by adjoining \(S\).  Thus its degree is five, proving the second
identity in \eqref{eq:base-invariants}.
\end{proof}

Since the determinants of \(A\) and \(L\) are \(1/2\) and \(-1\),
respectively, \eqref{eq:base-invariants} gives
\[
 \det D\Phi=(-1)(-2)\left(\frac12\right)=1,
 \qquad
 \gdeg(\Phi)=5.
 \tag{2.15}\label{eq:normalized-invariants}
\]

\subsection{Why the target line appears}

At this point the map is fixed.  What remains is to see the whole
intersective family inside its target.  Set \(\Pi=1\).  For \(a\in\Q\), put
\[
 \widetilde y_a=
 \left(1,\frac{32a}{9},\frac{8a+1}{3}\right).
 \tag{2.16}\label{eq:unnormalized-target}
\]
The elementary identity
\[
 \left(S-\frac12\right)^3\left(S^2+\frac34\right)
 =G(S)-\frac3{32}
 \tag{2.17}\label{eq:centered-identity}
\]
gives
\begin{align}
 E_{\widetilde y_a}(S)
 &=
 \left(\left(S-\frac12\right)^3-a\right)
 \left(S^2+\frac34\right).
 \tag{2.18}\label{eq:target-factorization}
\end{align}
After \(X=S-\tfrac12\), this becomes
\[
 (X^3-a)(X^2+X+1).
 \tag{2.19}\label{eq:translated-factorization}
\]
The target line is now visible.  Changing \(a\) changes the target, not the
map.  The centered quadratic has no linear term, so the parameter is
absorbed exactly by \(B\) and \(C\).

The target \(y_a\) in \eqref{eq:target-family} is \(L(\widetilde y_a)\).  The source
automorphism \(A\) and Proposition~\ref{prop:gauge} therefore give
\eqref{eq:main-fiber} whenever \eqref{eq:translated-factorization} is squarefree.

Its two factor discriminants and their resultant are
\[
 \operatorname{disc}(X^3-a)=-27a^2,\qquad
 \operatorname{disc}(X^2+X+1)=-3,\qquad
 \operatorname{Res}=(a-1)^2.
 \tag{2.20}\label{eq:disc-resultant}
\]
Thus every \(a>1\) in \(\mathcal A\) gives a reduced fiber of rank five.
By \eqref{eq:normalized-invariants}, the fiber is full.
